\documentclass{beamer}
%\usetheme{Madrid}
%\usetheme{Berlin}
%\usepackage{german}
%\usepackage[latin1]{inputenc}
%\usepackage[T1]{fontenc}
%\mode<presentation>
 %{
 \usetheme{Madrid}
 \setbeamercolor*{palette primary}{use=structure,fg=black,bg=p7408C}
 \setbeamercolor*{palette secondary}{use=structure,fg=black,bg=p7406U}
 \setbeamercolor*{palette tertiary}{use=structure,fg=black,bg=p716C}
 %}
 
 \definecolor{p7408C}{HTML}{F6BE00}
  \definecolor{p7406U}{HTML}{FCB315}
 \definecolor{pyelloy}{HTML}{F6C100}
 \definecolor{p716C}{HTML}{F17C0E}
 
  \definecolor{p8C}{HTML}{A19589}
 \setbeamercolor{itemize item}{fg=p8C} % all frames will have red bullets
 \setbeamertemplate{itemize item}[triangle]
 \setbeamertemplate{itemize subitem}{\color{p7408C}$\blacktriangleright$}

\setbeamercolor{block title}{bg=p7408C,fg=black}

\usepackage{etex}
\usepackage{tikz,pgfplots}
\usepackage{picture}

%\usepackage{german}
%\usepackage[latin1]{inputenc}
\usepackage[T1]{fontenc}
\usepackage{booktabs}
\usepackage{graphicx}
\usepackage{epsfig}
\usepackage{changebar}
\usepackage{amsmath}
\usepackage{mathtools}
\usepackage[T1]{fontenc}
\usepackage[ansinew]{inputenc}
\usepackage{setspace}
\usepackage{multicol}
\usepackage{multirow}
\usepackage{dcolumn}
\usepackage{color}
\usepackage{colortbl}
\definecolor{lightgray}{gray}{0.8}
\definecolor{lightblue}{rgb}{0.8,0.8,0.8}
\usepackage{hyperref}
\usepackage{xcolor} 
\usepackage{array} 
\usepackage{subfig}
\usepackage{makecell}
\usetikzlibrary{arrows,positioning} 


\graphicspath{{/Users/jupiter/Dropbox/elements/coala/graphs/}{../../../../Dropbox/elements/spatialrebels/graphs/}{../../../../Dropbox/elements/spatialrebels/graphspres/}}



\begin{document}

\title[UCL]
{
Cross-Validating Binary-Time-Series-Cross-Section Data}
\author[PSA PolMeth]
{
Gokhan Ciflikli \& Nils W. Metternich   \\[3ex]
\footnotesize{\emph{COnflict Analysis LAb (COALA)}}\\
\textbf{Funded by the Economic and Social Research Council (UK)}}
\institute[]
{		Department of International Relations\\
		London School of Economics and Political Science\\
		g.ciflikli@lse.ac.uk \\[4ex]
		
		Department of Political Science\\
		University College London\\
		n.metternich@ucl.ac.uk\\[4ex]
		
		}

\date{\today}

\begin{frame}
\titlepage
\end{frame}

\begin{frame}{Cross-validation}
\begin{itemize} 
\item Model validation technique to generalize to an independent data set
\item Done by partitioning datasets into training and test sets.
\end{itemize}
\begin{itemize}
\item Good cross-validation strategies are important for:
\begin{enumerate}
\item Prediction
\item Preventing overfitting
\item Robust inference
\end{enumerate}
\end{itemize}

\end{frame}

\begin{frame}{Binary-Time-Series-Cross-Section Data}
\begin{itemize}
\item Very common data format (Beck et al. +2500 citations, Carter and Signorino +1000 citations)
\item Binary outcome variable
\item Across observations 
\item Over time
\end{itemize}
\end{frame}

\begin{frame}{What is special about BTSC?}
\begin{itemize}
\item In some instances: Special case of duration data (Beck et al. 1998)
\item This is the case if underlying risk of occurrence varies over time
\item Even though this means that cases are IID, risk information can only be gathered by observing complete process
\end{itemize}
\end{frame}

\begin{frame}{When should you care about type of CV}
\begin{itemize}
\item If time doesn't matter you are fine
\begin{itemize}
\item if it does...
\end{itemize}
\item If time has a linear trend you are fine
\begin{itemize}
\item if it doesn't ...
\end{itemize}
\item then type of cross-validation matters
\end{itemize}
\end{frame}


\begin{frame}{Cross-validation approaches}
\begin{itemize}
\item Rolling-window (starting with about 60 percent of data)
\item 10-fold cross-validation
\item Group-wise cross-validation (10-fold)
\end{itemize}
\end{frame}

\begin{frame}{DGP}
\begin{itemize}
\item Survival data (Hendry, 2014)
\item IDs: 1000 
\item Mean Life: 17
\item Samples: 100
\end{itemize}
\centering
\includegraphics[scale=.3]{/Users/jupiter/Desktop/hist_cox.png}
\end{frame}


\begin{frame}{Logit: Y=X+t+t$^2$+t$^3$}
\centering
\includegraphics[scale=.3]{../../../../Desktop/results_cox.png}
\end{frame}


\begin{frame}{Results}
\centering
\begin{tabular}{ccc}
Cross-validation & True-Positive Rate & True-Negative Rate\\ \hline
Group-wise CV (10-fold) & 0.003416897 & 97.039709700 \\
10-fold CV & 0.002451740 & 97.146360996 \\
Rolling-window CV & 0.071861408 & 92.377115650  \\
\end{tabular}
\end{frame}



\begin{frame}{DGP}
\begin{itemize}
\item Binomial with t+t$^2$+t$^3$
\item Observations: 1000
\item Samples: 100
\item Similar to Carter and Signorino
\end{itemize}
\centering
\includegraphics[scale=.3]{../../../../Desktop/hist_logit.png}
\end{frame}


\begin{frame}{Logit: Y=X+t+t$^2$+t$^3$}
\centering
\includegraphics[scale=.3]{../../../../Desktop/results_logit.png}
\end{frame}

\begin{frame} {Logit}
\begin{tabular}{ccc}
Cross-validation & True-Positive Rate & True-Negative Rate\\ \hline
Group-wise CV (10-fold) & 25.20300 & 41.17200 \\
10-fold CV & 25.26900 & 41.18800  \\
Rolling-window CV & 34.38667 & 26.12000   \\
\end{tabular}
\end{frame}
 



\begin{frame}{Other preliminary insights}
\begin{itemize}
\item Smaller samples type of CV matters more
\item Especially if risk is increasing over time
\item And risk `peaks' at particular time points
\end{itemize}
\end{frame}





\begin{frame}{Conclusion and looking ahead}
\begin{itemize}
\item Settle on one DGP
\item Further Monte-Carlo experiments
\item Currently showing differences in existing work
\item Providing diagnostics for best CV type
\item Understanding CV in auto-regressive data (non-IID)
\end{itemize}
\end{frame}

\end{document}





